\section{The probe and its averaging scales}\label{gex:source-section}

\subsection{Field, characters and calibration}

Let $F=\Q(\sqrt{-3})$, $\omega=e^{2\pi i/3}$ and
$\lambda=1+2\omega$. Its ring of integers is $\OO=\mathbb Z[\omega]$.
For an element or ideal write $q_a=\N(a)$; for a nonzero element put
$\alpha(a)=a/|a|$. Good ideals are represented by their unique
primary generators, congruent to one modulo $3$.

Fix a primitive finite-order target character $\eta$. The fixed ray
group $\mathcal T$ presents the supplementary and reciprocity phases
$G$ and $\mathcal R$ of \cite[Section 4]{OAI26}; it is independent
of the target. Choose the finite excluded set $S$ to contain the
primes over $6$, the full defining-modulus supports of $\eta$ and
$\widehat{\mathcal T}$, and every prime below a fixed cutoff $P_0$.
The defining modulus includes redundant primes at which an induced
presentation is extended by zero. This entire datum is fixed before
the scale $Z$ and moving prime labels are chosen.

Let $I_\eta(X,Y,Z)$ be exactly the calibrated probe of
\cite[(6.1), Section 6.1]{OAI26}, with nonnegative nonzero annular
tests $W_0,W_1$ and Gaussian Mellin weight $e^{t^2}$.
We retain its auxiliary factors $b_*,\xi,\tau,\Xi,G,\mathcal R$
and normalized cubic Gauss sum $\gamma_2$.
Here $b_*$ generates the squarefree product of $S$;
$\xi\bmod b_*$ is primitive, of order dividing six and nonprincipal
at every prime factor, $|\tau|=1$, and
$\Xi(a)=\xi(a)\chi_a(b_*)$.
The sextic symbol $\chi_a(v)=(v/a)_6$ has its original zero on
nonunits. Thus a sixth power of the symbol is a coprimality
indicator. Inverses of $\Xi$ are taken only on units, and no
character is replaced by its primitive presentation.

The completed indices are $cn^3$, with $c$ squarefree and $n$
unrestricted; they may share primes. All ideal variables avoid $S$.
For a squarefree product $D$ of slot primes, define $I_{\eta;D}$
by inserting $\one_{D\mid cn^3}$ in this same completed row.
The auxiliary phases and the weights $q_c^{-1/2}q_n^{-1}$
remain unchanged. This convention fixes the source in both the
direct estimate and its Poisson representation.

\subsection{Compensation and geometry}

Take an even fixed slot count $K$, positive lengths $\ell_i$ with
$\sum_i\ell_i=\ell$, and nonnegative nonzero annular weights $W_i$.
Put $P_i=Z^{\ell_i}$. The underlying prime windows are disjoint
before ray or excluded-set restrictions. Set
\[
 \mathcal P_i(Z)=\{p\notin S:\ p\text{ prime},\ p\in1_{\mathcal T},
                 \ q_p/P_i\in\operatorname{supp}W_i\}.
\]
For a tuple $\boldsymbol p$ and subset
$J\subseteq\{1,\ldots,K\}$, put $p_J=\prod_{i\in J}p_i$.
Use the same finite operation as \cite[(12.5)]{OAI26}:
\begin{equation}\label{gex:compensated-source}
\begin{split}
 \mathcal I_{\eta,\mathrm{mod}}(Z)=
 \sum_{\boldsymbol p\in\prod_i\mathcal P_i(Z)}
 \prod_iW_i(q_{p_i}/P_i)
 \sum_J(-1)^{|J|}q_{p_J}^{-3/2}\overline{\eta(p_{J^c})}\,
 I_{\eta;p_{J^c}}\!\left(\frac{X}{q_{p_J}},\frac{Y}{q_{p_J}},
                           Zq_{p_{J^c}}\right).
\end{split}
\end{equation}
Every window remains at its original $P_i$. Membership in
$1_{\mathcal T}$ does not imply $\eta(p)=1$; the conjugate target
factor in \eqref{gex:compensated-source} is retained.

We choose
\begin{equation}\label{gex:geometry}
\begin{gathered}
 \varpi=\frac1{6000},\qquad b=\frac18,\qquad
 \ell=\frac16+\varpi,\qquad h=\frac{13}{16}+\frac{3\varpi}{2},\\
 l_x=\frac{17}{48}-\frac\varpi2,\qquad
 l_y=\frac{23}{48}-\frac\varpi2,\qquad M=\frac56-\varpi,\\
 X=Z^{l_x},\quad Y=Z^{l_y},\quad
 M+\ell=1,\quad l_y-l_x=b,\quad h=1-l_x+\ell,\\
 \theta=\frac78-\frac\varpi4=\frac{20999}{24000}.
\end{gathered}
\end{equation}
At $\varpi=0$ these are the values of \cite[(12.4)]{OAI26}; the scalar
perturbation $\varpi$ is the only change. The fine slot system is
chosen later from the moment mesh. It remains fixed as $Z$ tends to infinity.

For a rescaled subset put $d=\sum_{i\in J}\ell_i$. Write
$q_{p_J}=Z^d r_J$, with $r_J$ in one fixed positive compact
interval, uniformly over the tuples. Then the actual low scales are
\[
 X'=Z^{l_x-d}/r_J,\quad Y'=Z^{l_y-d}/r_J,\quad
 Q=q_{b_*}X'Y'\asymp Z^{M'},\qquad
 M'=1-\ell-2d,\quad \ell'=\ell-d.
\]
In particular $l_x-\ell,l_y-\ell$ and $M-2\ell$ are positive.
All rescaled row and lattice scales eventually exceed one, uniformly
in the fixed annular ratios.

\subsection{The shared Mellin representation}

Let $\widetilde W_0(|z|^2)$ be the planar Fourier transform for the
self-dual measure $(2/\sqrt3)\,dx\,dy$, and put
\[
 M(z)=\int_0^\infty\widetilde W_0(r)r^{z-1}\,dr,
 \qquad \widehat W_1(w)=\int_0^\infty W_1(r)r^{w-1}\,dr.
\]
Both transforms have rapid vertical decay in the relevant fixed
strips; $M$ is holomorphic for $\Rea z>0$ and $M(1/6)>0$
by \cite[Section 7.1]{OAI26}.

The exact Poisson identity of \cite[Section 16.1]{OAI26} gives
\begin{equation}\label{gex:high-identity}
\begin{split}
 \mathcal I_{\eta,\mathrm{mod}}(Z)
 =\frac1{(2\pi i)^3}\int_{(3)}\int_{(3)}\int_{(2)}
 &X^{1/2-z}Z^{s+z-1}Y^{w-1}
   e^{(s+z-1)^2}M(z)\widehat W_1(w)\\
 &\times\sum_u^{(6)}q_u^{-z}\overline{\xi(u)}
 \frac{\zeta_F^S(6z)L_F^S(w,\chi_\bullet(u))}{
       L_F^S(s,\eta\overline{\chi_\bullet(u)})}
 \mathfrak H_{\eta,u,Z}(s,w,z)\,dz\,dw\,ds.
\end{split}
\end{equation}
The sum includes all nonzero sixth-power-free element rows prime
to $S$, with unit factors. Superscript $S$ denotes deleted Euler
factors. The correction $\mathfrak H$ is the full quotient-free
selected tuple correction of \cite[(16.3)]{OAI26}. At a selected
prime it uses the holomorphic factor $G_p$; all unselected primes
retain $H_p$. This definition is valid also at zeros of individual
raw local quotients. Identity~\eqref{gex:high-identity} concerns
exactly the finite physical expression~\eqref{gex:compensated-source}.

At the principal double residue $u=1,w=1,z=1/6$, the outside
power is $X^{1/3}Z^{s-5/6}$ and the slot factor supplies
$Z^{\ell/6}$. Thus the principal exponent of \cite[(12.3)]{OAI26} persists:
\begin{equation}\label{gex:principal-exponent}
 C(s)=s-\frac56+\frac{l_x}{3}+\frac\ell6
     =\frac{l_x}{2}+s-1+\frac h6=s-\frac{11}{16}.
\end{equation}
Increasing $\ell$ therefore leaves $C(s)$ unchanged. The low
estimate, by contrast, decreases with $l_x$; this is the source
of the improvement.
